\section{The same original global terminal projection is biholomorphic off the true zero locus}

The off-zero clause of principalization in \cite{wlodarczyk05-hironaka}
is retained by the actual global terminal gluing.  The proof uses genuine
finite terminal projections and the already constructed original saturated
ranges, on a base which need not be compact.  It constructs an actual
biholomorphism of open complements without presupposing a global blowup
schedule.

\begin{theorem}[The actual original global projection is surjective]
\label{hir:original-global-terminal-projection-surjective}
\lean{Hironaka.HironakaOriginalGlobalOffZero.originalTerminalGluedProjection_surjective}
\leanok
\uses{hir:original-input-covering-local-canonical-terminal-atlas,hir:original-input-local-terminal-glued-whole-continuous-projection,hir:principalization-local}
For every original marked ideal in the invariant domain, its actual global
terminal projection $P:Y\to X$ is surjective, without compactness of $X$.
\end{theorem}
\begin{proof}
Fix $x\in X$.  Its internally chosen original atlas open $U_x$ contains
$x$.  The actual local terminal tower projection $p_x:Y_x\to U_x$ is
surjective, since each genuine blowup is surjective and the finite composite
is surjective.  Choose a point above $x$ in this actual tower and apply
the original local inclusion.  The whole projection dictionary identifies
its global image with $x$.  This constructs a preimage of every original
base point.
\end{proof}

\begin{theorem}[The actual original global projection is an isomorphism off the original ideal]
\label{hir:original-global-terminal-isomorphism-off-true-ideal}
\lean{Hironaka.HironakaOriginalGlobalOffZero.originalLocalTerminalProjection_isIsomorphismOffZero,Hironaka.HironakaOriginalGlobalOffZero.originalTerminalGluedProjection_isIsomorphismOffZero}
\leanok
\source{wlodarczyk05-hironaka,bierstone-milman08-functoriality}
\uses{hir:original-global-terminal-projection-surjective,hir:original-global-principalization-true-support,hir:original-terminal-range-true-principalized-ideal,hir:original-input-terminal-glued-local-range-whole-open-preimage}
For an original marked ideal $M$ in the invariant domain, with whole
ambient submanifold and mark one, the same original global projection
$P$ maps $P^{-1}(X\setminus Z(M.\mathfrak a))$ bijectively to
$X\setminus Z(M.\mathfrak a)$, and its differential is invertible at
every point of this inverse image.
\end{theorem}
\begin{proof}
On each actual local terminal tower, mark-one principalization identifies
the total transform with the tower's effective exceptional divisor.
Its support is the inverse image of the original zero locus.  Every
exceptional point of the finite tower lies in that support.  Genuine
finite blowup geometry therefore says that the tower projection is
bijective and has invertible differential off that true zero locus.

The global projection is surjective by the preceding theorem, so it is
surjective over the complement.  Suppose two global points in this
complement have the same original projection.  Represent the first in
one actual local terminal chart.  Saturation of that range puts the
second point in the very same local range; take its actual representative
there.  The whole projection dictionary gives equal images under the
same local tower projection, in its off-zero region.  Its proved
injectivity identifies the representatives, and then the global points.

For invertibility of the differential, represent an arbitrary off-zero
global point by a genuine local point.  The composite of the local tower
projection with the original base-open inclusion is unramified there:
the tower differential is invertible, as is the base-open inclusion's
differential.  This entire composite equals the global projection
composed with the actual holomorphic local inclusion.  The chain rule in
equal complex dimensions implies that an invertible composite
differential has invertible global-projection differential.  This proves
the claimed nonexceptionality at every off-zero global point.
\end{proof}

\begin{theorem}[The true global open complements are genuinely biholomorphic]
\label{hir:original-global-terminal-genuine-off-zero-biholomorphism}
\lean{Hironaka.HironakaOriginalGlobalOffZero.originalTerminalOffZeroBaseOpen,Hironaka.HironakaOriginalGlobalOffZero.originalTerminalOffZeroSourceOpen,Hironaka.HironakaOriginalGlobalOffZero.originalTerminalOffZeroBiholomorphism,Hironaka.HironakaOriginalGlobalOffZero.originalTerminalOffZeroBiholomorphism_point}
\leanok
\source{wlodarczyk05-hironaka,bierstone-milman08-functoriality}
\uses{hir:original-global-terminal-isomorphism-off-true-ideal,hir:original-input-global-terminal-whole-projection-holomorphic}
The actual opens $X\setminus Z(M.\mathfrak a)$ and its true inverse image
inside the original global terminal manifold are genuinely biholomorphic.
The forward map is the restriction of the same original global projection
$P$, with both holomorphic directions proved.
\end{theorem}
\begin{proof}
Coherence makes the original zero locus closed.  Its complement is
therefore an actual open submanifold, and holomorphicity of $P$ makes
its true inverse image open.  The preceding theorem constructs a
bijection between these actual open submanifolds.  Its forward map is
holomorphic because it is the original projection restricted to these
opens.  Fix a point of the target complement and its actual inverse
representative.  The original projection is unramified at that
representative, so the holomorphic inverse function theorem supplies
a local holomorphic inverse on a neighborhood of the target point.
Intersect that neighborhood with the complement.  Every local inverse
value lies in the true source complement because its projection is the
target point.  Global off-zero injectivity therefore identifies this
local inverse with the actual inverse of the constructed bijection on
the whole neighborhood.  The actual inverse is holomorphic there, and
this holds at every target point.  Bundle the bijection and the two
proved holomorphic directions into the genuine biholomorphism.
\end{proof}
